% !TeX root = ../../protocolo.tex

In the the past section, the optical response of AuNPs metasurface was tudied considering the contributions of the substrate effect under realistic conditions. A partial embedding along enhanced by the multiple scattering within the AuNPs ensemble was proposed as the main mechanism of experimental spectral shift of the metasurfaces's response relative to the theoretical results, which consider an ideal scenario. To continue this line of investigation, a new set of metasurfaces were fabricated as explained above but the annealing temperature was varied, so a controlled average partial embedding of the metasurface can be achieved due to the fusion  melting of the substrate supporting the AuNPs. In this section, the first optical characterization of the fabricated samples are presented, however it is stressed that SEM and Atomic Force Microscopy (AFM) images are needed to characterize the size distribution the AuNPs and to  estimate the partial embedding of the metasruface. 

In Fig. \ref{fig:extexp} it is shown the transmissivity spectra of the synthesized  AuNPs metasurfaces fabricated by thermal annealing continuous gold films on Corning glass of different nominal thicknesses $d$. The thermal treatment was performed at two distinct temperatures ---$500^\circ$ [Fig. \sfref{fig:extexp}{a)}] and $550^\circ$  [Fig. \sfref{fig:extexp}{b)}]--- on gold films varying from $3$ nm of thickness to $15$ nm for the sample annealed at $500^\circ$ nm and  from $3$ nm to $11$ nm for the $550^\circ$ annealed sample. The color code corresponds to one value of the initial thicknesses of the samples with the exception of the red curve, which corresponds to the transmissivity of a clean substrate employed as a reference.

\begin{figure}[h]\centering
	\includegraphics[scale = 1]{exp/Fig-Extinction.png}
	\caption{%
		Transmissivity spectra of the synthesized metasurfaces  after thermal annealing at a \textbf{a)} $500^\circ$ and \textbf{b)} $550^\circ$ temperature at normal incidence ($\theta_i = 0^\circ$) as function of the incidence wavelength. The color legend on the right indicates the nominal thicknesses of the initial deposited Au films prior to annealing, while he red curve corresponds to the reference obtained by measuring the transmissivity of a clean substrate. }
	\label{fig:extexp}
\end{figure}

For either of the analyzed samples sets, that is considering a fixed annealing temperature, the extinction spectra minima exhibit a redshift that increases as the initial thickness of the film does  up to the value of  $9$ nm and, similarly, the overall value of the transmissivity within the visible spectra decreases for the thicker samples. This results are consistent with the fabricated samples addressed in Section  \ref{s:substratep} where SEM images were available, thus larger AuNPs are expected in the metasurface for thicker films. It is noteworthy that for the samples synthesized with $d\leq 11$~nm the transmissivity minima are spectrally located in the visible range at wavelengths $\lambda < 700$~nm while the samples with an initial thickness of $d = 13$~nm and $d=15$~nm heated at $500^\circ$C ---yellowish lines in Fig. \sfref{fig:extexp}{a)}--- present the broadest and most redshifted minima of all samples, which  suggest the agglomeration of the AuNPs or the synthesized of non-spherical particles. Hereby the discussion is centered on the metasurfaces fabricated with an initial thickness of $d\leq 11$~nm.

 The spectral behavior of the samples become more pronounced for the highest annealing temperature, indicating greater extinction from the  AuNP ensembles on the substrate and a optical response redshifted with increasing thermal treatment, both effects identified theoretically for partially embedded AuNPs, which can be obtained by a softening process of the  Corning glass substrate  for higher temperatures. 

The experimental reflection spectra  of the synthesized metasurfaces is shown in Fig. \ref{fig:atrexp}. The metasurfaces were illumintaed in an ATR configuration at an angle of incidence $\theta_i = 45^\circ$ at

Figure \ref{fig:ATRThick} displays the reflectivity spectra of annealed gold thin films with different initial thicknesses, measured at an incident angle of $45^\circ$ for both $p$- and $s$-polarized light. The measurements were performed after thermal annealing at 500°C (panel a, top) and 550°C (panels a-bottom and b). The influence of film thickness and polarization on the position and depth of the reflectivity minima is clearly observed. In all cases, increasing the film thickness results in a shift and broadening of the plasmonic resonance dip. The shaded region highlights the wavelength range in which the minima of most reflectivity curves are located, emphasizing the consistent spectral region of plasmonic response regardless of polarization or temperature.	


\begin{figure}[h]\centering
	\includegraphics[scale = 1]{exp/Fig-ATRThick.png}
	\caption{%
		Reflectivity spectra of thermally annealed Au thin films with varying initial thicknesses, measured at an incident angle $\theta_i = 45^\circ$ for (a) $p$-polarized light at 500°C (top) and 550°C (bottom), and (b) $s$-polarized light at 550°C. The shaded area marks the spectral region containing the reflectivity minima for most film thicknesses. A clear dependence on thickness, polarization, and annealing temperature is observed, affecting both the position and depth of the plasmonic resonance.}
	\label{fig:atrexp}
\end{figure}